%===============================================================================
%
%         FILE: cv-sv.tex
%
%        USAGE: ---
%
%  DESCRIPTION: ---
%
%      OPTIONS: ---
% REQUIREMENTS: ---
%         BUGS: ---
%        NOTES: ---
%  ORIG AUTHOR: Samuel Gabrielsson <samuel.gabrielsson@gmail.com>
% ORGANIZATION: ---
%      VERSION: 1.0
%      CREATED: 2026-04-31 13:46:31
%     REVISION: ---
%      CHANGES: ---
%
%===============================================================================
\documentclass[10pt,a4paper,sans]{moderncv}

% Configure
\moderncvstyle{classic}             % Style options: casual, classic, banking, oldstyle and fancy
\moderncvcolor{blue}                % Color options
\usepackage[scale=0.75]{geometry}   % Adjust the page margins

%\setlength{\hintscolumnwidth}{0.25\textwidth}

% Personal data
\name{Samuel}{Gabrielsson}
%\title{Resumé Title}
\address{Saitama AB}{Lindhagensgatan 59}{112 43 Stockholm}
\phone[mobile]{+46~(0)72~152~42~28}
\email{samuel.gabrielsson@gmail.com}
%\homepage{www.ludd.ltu.se/~proximus}

\photo[64pt][0.4pt]{profile-picture}

% Social icons
\social[linkedin]{samuelgabrielsson}
\social[github]{proximus}

% BibTex adjustments
%   - Show numerical labels in the bibliography
\renewcommand*{\bibliographyitemlabel}{[\arabic{enumiv}]}

\begin{document}
%===============================================================================
% Resume/CV Content
%===============================================================================
\makecvtitle

%===============================================================================
% Personal Profile
%===============================================================================
\section{Om Mig}
Med över 15 års erfarenhet inom omr{\aa}den som \textbf{Senior Embedded Linux
and Kernel Developer}, \textbf{Scrum Master}, \textbf{Produktägare} samt
specialist i \textbf{Systemintegration}, \textbf{IT-Automation},
\textbf{Testramverk-Arkitektur} och \textbf{Testfall}, kombinerar jag teknisk
expertis med starkt ledarskap.
\newline
\newline
Ett tekniskt bidrag jag är särskilt stolt över är mitt arbete med Linux-kärnan,
där jag utvecklade patchar för ECC memory scrubbing – en funktion som skyddar
mot strålningsinducerade bitfel och bland annat används i system som
\textbf{NASA:s Ingenuity Mars-helikopter}, som utförde den första kontrollerade
motorflygningen på en annan planet.
\newline
\newline
Jag är en problemlösare som närmar mig varje utmaning med nyfikenhet och
uth{\aa}llighet, och jag värdesätter samarbete och en positiv arbetsmiljö. Genom
pålitlighet, öppenhet och kunskapsdelning bygger jag förtroende hos både
kollegor och kunder.

%===============================================================================
% Work experience
%===============================================================================
\section{Erfarenhet}
% 20240314 - 20260227
\cventry{Mar 14, 2024-- Mar 01, 2026 }{Utvecklare, Scrum Master \& Produktägare}{Saitama AB}{Stockholm}{Sweden}{
\begin{itemize}
\item \textbf{\textit{Konsult på Ericsson AB - Radio Software}}\newline{}
Ett tv{\aa}{\aa}rigt kontrakt inom RadioSW RX/Uplink-teamet, där jag
balanserade tre roller samtidigt: utvecklare, Scrum Master och Produktägare för
Massive MIMO Granada-programmet samt BALI-projektet för radioprodukten AAS AIR
6494.
    \begin{itemize}
    \item Utvecklare med fokus på Uplink/Rx-delen av mjukvaran, bestående av drivrutiner och testfall.
    \item Driva arbetssättet genom dagliga stand-ups, sprintgranskningar, retrospektiv, refinement och sprintplanering.
    \item Produktägare för Uplink/Rx-delen av produkten. Initierade backloggen för hela RadioSW BALI-projektet i Jira. Skapa, förfina och planera backloggen för teamet.
    \end{itemize}
    Utvecklingsverktyg: C/C++, Java Yocto, Git, Gerrit, Jenkins, Jira och NeoVim.
\end{itemize}
}
\bigskip

% 20230601 - 20231231
\cventry{Jun 01, 2023-- Dec 31, 2023 }{Embedded Developer \& Scrum Master}{Saitama AB}{Stockholm}{Sweden}{
\begin{itemize}
\item \textbf{\textit{Konsult på CrossControl AB - Uppsala}}\newline{}
Scrum master och Embedded Linux-utvecklare som ansvarade för ett Windows-team
och ett Linux-team. Utvecklade drivrutiner och applikationer för produkter som
CCPilot X1200 och X900 som är en 12-tums display med prestanda i nivå med en
bärbar dator för industriella fordon, utformad för att uppfylla krävande
miljökrav. Den stöder moderna och nästa generations verktyg som sömlöst kan
integrera flera kameraflöden i ditt HMI-system, vilket förbättrar operatörens
sikt och säkerhet.
    \begin{itemize}
    \item Leda teammöten, dagliga stand-ups, sprintgranskningar, retrospektiv, refinement och sprintplanering. Hålla backloggen uppdaterad med tillräcklig kvalitet. Skapa väl dokumenterade uppgifter och fördela dem med ett bra arbetsflöde.
    \item Utveckla Linux-kärndrivrutiner och intern mjukvara för att stödja ljussensorn, EEPROM, automatisk bakgrundsbelysning, knappsatser, audio codec, PXE bootable installationsavbilder för kund och produktion, I2C-kommunikation samt CAN-FD MCU.
    \end{itemize}
    Utvecklingsverktyg: C/C++, Python, Bash, Yocto, Devtool, Jira, Git, BitBucket, Jenkins och NeoVim.
\end{itemize}
}
\bigskip

% 20221001 - 20230331
\cventry{Okt 01, 2022-- Mar 31, 2023 }{Digital Radio Algorithms Developer}{Saitama AB}{Stockholm}{Sweden}{
\begin{itemize}
\item \textbf{\textit{Konsult på Ericsson AB - Hårdvaruavdelningen}}\newline{}
Pre-development av Ericssons Shared Digital Pre-Distortion (DPD). Arbetade
tillsammans med labbet i Ottawa för att effektivt linearisera effektförstärkare
(PA:er) i 5G-radioenheter. I stället för att använda en separat DPD för varje
antenngren gör Shared-DPD det möjligt för en enda DPD att hantera flera grenar
eller grupper av effektförstärkare med hjälp av en intern switch. Detta
minimerar olinjära beteenden, reducerar distortion inom bandet och minskar
emissioner utanför bandet, vilket maximerar energieffektiviteten i
5G-basstationer. Ericssons kommande 6G-radioprodukter kommer att använda
Shared-DPD.
    \begin{itemize}
    \item Utveckla Digital Pre-Distortion (DPD)-mjukvara i Matlab och C för att öka effektiviteten hos effektförstärkare.
    \item Använda Xilinx SDK eller Vitis IDE för att implementera och felsöka FPGA samt mjukvaruladdning på målplattformen.
    \end{itemize}
    Utvecklingsverktyg: C, MATLAB, Yocto, Git, Xilinx SDK, GDB och Vim.
\end{itemize}
}
\bigskip

% 20181001 - 20220930
\cventry{Okt 01, 2018-- Sep 30, 2022 }{Senior Research Software Engineer}{Saitama AB}{Stockholm}{Sweden}{
\begin{itemize}
\item \textbf{\textit{Konsult på Ericsson AB - Systems and Technology (R\&D)}}\newline{}
Plattformsforskning och utveckling av 5G- och 6G-testbäddar, från fältförsök
till prototyper. Ansvarade för basbandsapplikationer på högre lager som körs på
ett realtidsoperativsystem, med fokus på att minimera CPU-jitter och
NIC-latens.
    \begin{itemize}
    \item Utveckla en UDP-server och aktivera hårdvarutidsstämpling på NIC för att mäta paketens round trip time (RTT).
    \item Utveckla en högpresterande AF\_XDP-klient för paketbearbetning med zero-copy, för körning på Mellanox CX-6 200Gb/s och CX-7 400 Gb/s Ethernet nätverkskort som stöder XDP.
    \item Utvecklade applikation som körs på Ericssons Many-Core Architecture (EMCA), vilket är en avancerad ASIC med digital logik, mixed-signal-innehåll och processorer.
    \item Utveckla ett testramverk och testfall i Python, Pexpect och YAML.
    \item Skapa och konfiguera realtids-OS baserat på Ubuntu. Stödja avdelningen och se till att realtidsapplikationen för basbandet kan köra på OS:et med höga krav på låg CPU-jitter.
    \item Installera och ansvara för 5G och 6G hårdvaran och delar av mjukvaran i labbet.
    \item Implementera zero-touch-automation med Ansible på Ericssons egenutvecklade hårdvara samt standardhårdvara (COTS) som Supermicro- och Dell PowerEdge-servrar med BMC/IPMI/iDRAC samt Mellanox SN4700 400GbE-switchar med Cumulus. Automationen genomför fullständig OS-distribution på flera noder parallellt och konfigurerar BIOS, LDAP och PTP, installerar Linux eller Cumulus OS obevakat samt sätter upp realtidsegenskaper med låg CPU-jitter och låg NIC-latens.
    \end{itemize}
    Utvecklingsverktyg: C/C++, Assembler, MATLAB, Python, Ansible, YAML, Bash, Git, Jenkins and Vim.
\end{itemize}
}
\bigskip

% 20161201 - 20180930
\cventry{Dec 01, 2016-- Sep 30, 2018 }{Senior Software Architect \& Designer}{Realtime Embedded AB}{Stockholm}{Sweden}{
Fast anställd som konsult på Realtime Embedded med uppdrag hos olika kunder.
\begin{itemize}
\item \textbf{\textit{Konsult på NorthStar}}\newline{}
Utveckling av NorthStar Advanced Connected Energy (ACE) gateway för smarta
batterisystem som kommunicerar med varandra genom Bluetooth.
    \begin{itemize}
    \item Använda Python för att implementera extern larmhantering för batterisystemet. Säkerställa att mjukvaran kan trigga och styra reläer via GPIO-pinnar i gatewayen när korrekt signalering via IPC-meddelanden tas emot.
    \item Implementera funktionella tester samt bibliotek/paket med hjälp av 9pm-testframeworket.
    \end{itemize}
    Utvecklingsverktyg: Python, Tcl/Expect, Bash, Git, Gerrit, Jenkins och Vim.
\item \textbf{\textit{Konsult på Vanderbilt (f.d. Bewator)}}\newline{}
Utveckling av Linux-kärndrivrutiner för passerkontrollsystemet Omnis och
produkten E100 centralenhet som baseras på AT91SAM9x5-EK MCU.
    \begin{itemize}
    \item Skriva en GPIO-testapplikation för att styra reläer.
    \item Skriva en testapplikation för Atmel AT91 ADC för att läsa råa spänningsvärden.
    \item Skriva Linux-kärnpatchar och uppdatera device trädet för att aktivera USB-B gadget och Ethernet över USB.
    \item Skriva Linux-kärndrivrutiner och applikation för 74VHC595 8-bitars skiftregister med utgångslatchar för att via sysfs styra lysdioder, inportar och RS485-kortläsare.
    \item Skriva Linux-kärndrivrutiner och testapplikation för kortläsare med Magstripe (Clock/Data) och Wiegand (Data/Data).
    \item Skriva ett firmwareuppgraderings-/flashskript för att automatiskt uppdatera systemet via SD/MMC-kort vid uppstart.
    \item Modifiera och patcha bootloadern från Atmel, kallad at91bootstrap, för att åtgärda buggar och möjliggöra uppstart från SD/MMC eller NAND-flash.
    \end{itemize}
    Utvecklingsverktyg: C, Bash, Buildroot, Git och Vim.
\end{itemize}
}
\bigskip

% 20110610 - 20161123
\cventry{Jun 10, 2011-- Nov 23, 2016 }{Software Developer \& Scrum Master}{Ericsson AB}{Stockholm}{Sweden}{
Anställd på Ericsson i Kista som mjukvaruutvecklare och Scrum Master till Radio
Base Station 6000-serien. Utvecklade drivrutiner och simulatorn för RadioSW
teamen. För att spara tid och pengar så brukar vi emulera hårdvaran först och
därefter utveckla drivrutinerna tills de fysiska korten anländer från fabriken.
    \begin{itemize}
    \item Utveckling av radiomjukvarudrivrutiner för Radio SW-plattformar. Kortuppstart (board bring-up) och integration av hårdvaruprodukter såsom Antenna Integrated Radio (AIR), mikroradioenheter, nästa generations radioenheter och 5G-programmet.
    \item Skapa programmet Nova för att utvecklaren lättare kan köra simulatorn. Nova används flitigt även idag för att utveckla 5G och 6G radios på avdelningen.
    \item Modularisera delar av Radio SW-koden.
    \item Jobba i Ottawa, Kanada på korttidsuppdrag (3 månader) för hårdvaruavdelningen som mjukvaruintegratör för Platform 5.2 AIR 32-produkter efter köpet av Nortel.
    \item Introducera verktyg som Git, Gerrit, Jenkins, BitBake och Autotools till hela RadioSW och fasa ut ClearCase.
    \item Agera teknisk Scrum/Kanban Master för att leda och coacha ett tvärfunktionellt drivrutinsteam på 10 personer. Genomföra dagliga scrummöten, sprintretrospektiv, följa upp flödesstatistik, eliminera slöseri och stötta produktägaren.
    \end{itemize}
    Utvecklingsverktyg: C/C++, Green Hills Hardware Debugger Probe, Ruby, Tcl/Ecpect, Python, Bash, Git, Gerrit, Jenkins och Vim.
}
\bigskip

% 20080825 - 20110609
\cventry{Aug 25, 2008-- Jun 09, 2011 }{Software Developer \& Testare}{SAAB Combitech}{Stockholm}{Sweden}{
Fast anställd som konsult på SAAB Combitech på uppdrag hos Ericsson AB Linux Center i Älvsjö och Telefonplan. 
\begin{itemize}
\item \textbf{\textit{Konsult på Ericsson AB - Linux Development Center}}\newline{}
Utvecklade automatiserad testinfrastruktur och Linux-plattformsprogramvara för
inbyggda telekomsystem, inklusive LOTC, ENUX och CPBS. Arbetade genom hela
stacken, från testautomation till Linux-kärnan och plattformsutveckling.
Ericssons ENUX Linux-distribution användes för virtualiserad emulering av APZ
Classic-runtime-miljön för PLEX-baserade system såsom AXE.
    \begin{itemize}
    \item Arkitekterade och utvecklade ett Perl-baserat ramverk för automatiserad testning av Linux-validering i inbyggda system.
    \item Designade och implementerade testsviter och testfall i Tcl/Expect, Perl, Bash och Python, med integrerat stöd för NETCONF och YAML.
    \item Utvecklade test- och valideringssviter för LOTC, Xen-baserad SLES Virtual Platform (SVP) samt plattformsprodukter inklusive GSDI och IS på GEP-hårdvara.
    \item Utvecklade Linux-kernelpatchar för ECC (Error Detection and Correction) minnessanering (memory scrubbing), vilket förhindrade datakorruption och systemkrascher orsakade av strålningsinducerade bitfel (t.ex. kosmisk strålning och elektriskt brus).
    \item Bidrog med patchar till Linux-kärnan upstream.
    \item Gav teknisk support för Open Multimedia Platform (OMP) genom att felsöka och åtgärda defekter i produktpaketet.
    \end{itemize}
    Utvecklingsverktyg: C, Python, Perl, Tcl/Expect, YAML, Bash, Git, Gerrit and Vim. 
\end{itemize}
}
\bigskip

\cventry{2006--2007}{Research Assistant}{University of Toronto}{Toronto}{Canada}{
Arbetade som forskningsassistent för professor J. Christopher Beck vid
Department of Mechanical and Industrial Engineering, University of Toronto.
Tjänsten finansierades genom ett universitetsstipendium samt Toronto
Intelligent Decision Engineering Laboratory (TIDEL). Arbetet fokuserade på
operationsanalys, optimering och distribuerade metoder för att lösa NP-svåra
kombinatoriska problem.
\begin{itemize}
\item \textbf{\textit{Forskningsbidrag}}
    \begin{itemize}
        \item Utvecklade metoder inom artificiell intelligens och planering för modellering och syntes av säkerhetsprotokoll med hjälp av AI-planneringsramverk.
        \item Bedrev forskning inom optimering, parallella och distribuerade algoritmer samt kommunikationsprotokoll för lösning av kombinatoriska problem.
    \end{itemize}
\item \textbf{\textit{System- och algoritmutveckling}}
    \begin{itemize}
        \item Implementerade en distribuerad Tabu Search-algoritm i MPI för klusterbaserad beräkning och experimentell utvärdering.
        \item Designade och driftsatte en MPI-baserad klustermiljö i laboratoriet för forskning och testning av parallella algoritmer.
    \end{itemize}
\item \textbf{\textit{Publikationer}}
    \begin{itemize}
        \item Modelling Security Protocol Synthesis with AI Planning, University of Toronto / TIDEL.
        \item Solving Combinatorial Problems with Parallel Cooperative Solvers, Ninth International Workshop on Distributed Constraint Reasoning, Brown University, Providence, USA.
    \end{itemize}
\end{itemize}
Utvecklingsverktyg: MPI, C/C++, Python, Git and Vim.
}
\bigskip

%===============================================================================
% Misc activities/jobs
%===============================================================================
\subsection{Övrigt}
\cventry{2005--2006}{Styrelsemedlem}{Lule\aa \ University Academic Computer Society}{City}{Lule\aa}
{Ansvarade för 600 registrerade medlemmar.}
\cventry{2004--2005}{Kassör}{Lule\aa \ University Academic Computer Society}{City}{Lule\aa}
{Ansvarade för ekonomin med en budget på över 180 000 SEK. Hanterade inkomster och utgifter, skötte bokföringen, betalade fakturor och upprättade årsredovisning}
\bigskip

%\cventry{2000--2001, Summers 2002--2005}{Mail Carrier}{City-Mail}{Stockholm}{Sweden}{
%    \begin{itemize}
%    \item Sorting and distribution.
%    \item Updating the computer-based estate register.
%    \end{itemize}
%}
%\bigskip
%
%\cventry{1999--2000}{Substitute Teacher}{S\"odert\"alje Municipality}{S\"odert\"alje}{Sweden}
%{Teaching junior high school students math, chemistry and physics.}
%\bigskip
%
\cventry{1999--1999}{Militärtjänst, KA1 Vaxholm}{Svenska staten}{Stockholm}{Sweden}{
    \begin{itemize}
    \item Svenska marinen - Batteriledningsgruppbefäl vars uppgift är att handha den tekniska utrustning som används för det avancerade ledningen.
    \item Högvakt på kungliga slottet.
    \end{itemize}
}

%===============================================================================
% Education
%===============================================================================
\section{Utbildning}
\cventry{2001--2006}{Civilingenjör i Datateknik}{Luleå Tekniska Universitet}{Luleå}{Sweden}
{Med inriktning: Tillämpad Matematik}
\subsection{Examensarbete}
\cvitem{Titel}{\emph{A Parallel Tabu Search Algorithm for the Quadratic Assignment Problem}}
\cvitem{Handledare}{Professor J. Christopher Beck - University of Toronto \& Professor Inge Söderkvist - Luleå University of Technology}
\cvitem{Beskrivning}{Find optimal solutions for a quadratic assignment problem using the tabu search algorithm on a cluster.}
\bigskip

\cventry{1995--1998}{Naturvetenskapsprogrammet}{V\"asterg\aa rd High School}{S\"odert\"alje}{Sweden}{}

%===============================================================================
%  Skills
%===============================================================================
%\clearpage
\section{Färdigheter}
%\cvskilllegend*[1em]{}  % Adjust post spacing
\cvskillhead[-0.1em]    % Insert standard legend in English and adjust padding

% Programming
\cvskillentry*{Programming Languages:}
               {3}{Assembler}{1}{Återkommande arbete inom olika embedded-uppdrag.}
\cvskillentry{}{5}{C/C++}{20}{Jag använder C/C++ i de flesta stora Linuxbaserade inbyggda projekt, Linuxapplikationer, kärnmoduler och patchar.}
%\cvskillentry{}{2}{CSS}{2}{Used to develop my personal homepage a long time ago.}
\cvskillentry{}{4}{Haskell}{2}{Tillägnade mig rekursiv programmering under universitetsstudier.}
%\cvskillentry{}{3}{HTML}{2}{Used to develop my personal homepage a long time ago.}
\cvskillentry{}{4}{Java}{5}{Genom universitetsstudier och yrkeserfarenhet.}
%\cvskillentry{}{5}{\LaTeX}{10}{I've used \LaTeX{} in everything from thesis report, research papers, my company annual reports to this CV.}
\cvskillentry{}{4}{Lisp/Scheme}{3}{Har genomgått SICP-kursen vid universitetet och använt språken för att konfigurera min Emacs-editor.}
%\cvskillentry{}{4}{Maple}{3}{Used in mathematics courses at the university.}
\cvskillentry{}{4}{Matlab}{5}{Använt i matematikkurser vid universitetet samt under uppdrag och arbete inom Ericssons forskningsavdelning.}
\cvskillentry{}{4}{Pascal}{4}{Mitt första programmeringsspråk. Självlärd i Turbo Pascal sedan barndomen.}
\cvskillentry{}{5}{Perl}{7}{Utvecklade ett testramverk på Ericsson Linux DC som kunde köra testfall skrivna i Perl och Expect.}
%\cvskillentry{}{3}{PHP}{2}{Used to develop my personal homepage a long time ago.}
\cvskillentry{}{5}{Python}{15}{Utvecklade ett testramverk och tillhörande testfall. Python har även använts för diverse skriptning i större projekt genom åren.}
\cvskillentry{}{4}{Ruby}{2}{Skrivit testfall inom Radio SW. Mycket smidigt och användarvänligt skriptspråk med många praktiska fördefinierade funktioner.}
\cvskillentry{}{5}{Shell scripting}{23}{Använder Bash, sed, awk och reguljära uttryck till nästan alla mina arbetsprojekt samt hemma till t.ex. hobbyprojekt.}
\cvskillentry{}{5}{Tcl/Expect}{13}{Använt för automation och testfall i olika projekt.}

% Software Tools
\cvskillentry*{DevOps/Tools:}
               {5}{Ansible}{5}{Använt för zero-touch-automation tillsammans med Jinja2, YAML och Redfish.}
\cvskillentry{}{5}{Buildroot}{3}{Byggsystem för att bygga en hel Linux distribution.}
%\cvskillentry{}{3}{ClearCase}{3}{Version control system used at Ericsson but replaced by Git.}
\cvskillentry{}{4}{Docker}{3}{Använder hemma eller när jag behöver isolera applikationer hos kunden.}
\cvskillentry{}{5}{Gerrit}{10}{Verktyg för kodgranskning som använts hos Ericsson tillsammans med Git och Jenkins}
\cvskillentry{}{5}{Git}{15}{Använder det konstant hemma och på alla uppdrag.}
\cvskillentry{}{5}{Jenkins}{10}{Tillsammans med Gerrit och Git ger det ett mycket effektivt verktyg för CI.}
\cvskillentry{}{3}{SVN}{1}{Ett versionshanteringssystem som jag använde på ett uppdrag.}
\cvskillentry{}{5}{Vim/NeoVim}{20}{Mitt huvudsakliga textredigeringsverktyg för programmering och redigering.}
\cvskillentry{}{5}{Yocto/BitBake/Poky}{10}{Använt för att bygga inbyggda Linux-operativsystem på många uppdrag.}

% Operating Systems
%\cvskillentry*{OS:}{5}{Linux}{23}{Började med Red Hat, Slackware och Linux From Scratch, och övergick senare till Ubuntu som huvudsakligt system.}
%\cvskillentry{}{3}{OSE}{5}{The main OS in the old Ericsson Radio product.}
%\cvskillentry{}{4}{Solaris}{5}{Primary OS at the university.}
%\cvskillentry{}{5}{Windows/MS-DOS}{25}{Dual booting Windows only when playing PC games at rare occasions.}

\cvskillentry*[1em]{Methods}{5}{Scrum/Kanban}{10}{Mångårig erfarenhet från teamledarskap och agila arbetssätt.}

\section{Språk}
\cvitemwithcomment{Svenska}{Modersmål}{}
\cvitemwithcomment{Engelska}{Flytande}{}
%\cvitemwithcomment{Aramaic}{Fluent}{Ancient christian language spoken by my parents}
\cvitemwithcomment{Tyska}{Grundläggande}{}

%===============================================================================
% Reference
%===============================================================================
\section{Referenser}
\cvitem{}{Lämnas på begäran.}

%===============================================================================
% Publications
%===============================================================================
\nocite{*}
\bibliographystyle{plain}
\bibliography{publications} % publications is the name of a BibTeX file

\end{document}
